% Owner: boundary_d.  Independent proof; reserved by root as B.4.11.
% This is an auxiliary contracted surface, not the requested S^6 filling.
\begin{proposition}[B.4.11 (Proved): globalization at the contracted parabolic elliptic axis]
Let $S=S_8(a,\kappa)$ be the central parabolic surface in B.4.10,
let $E\subset S$ be its elliptic axis, and put
$N=N_{E/S}$, so $\deg N=-8$.
The contraction $\pi:S\to Y$ of $E$ is a compact normal surface
with a single isolated singular point $p$.
Its germ at $p$ is the vertex of the affine cone over
$(E,N^{-1})$.  In particular,
\[
 H^2(Y,\Theta_Y)=0.
\]
The restriction morphism from a semiuniversal deformation of $Y$
to a semiuniversal deformation of $(Y,p)$ is regular.
Consequently every prescribed analytic one-parameter deformation
of this isolated singularity is induced, after shrinking its
parameter disc, by a proper deformation of the compact surface $Y$.
If the prescribed local deformation smooths the singularity, the
global deformation has smooth fibres away from the central value.

This concerns globalization of a specified local deformation of
an auxiliary compact surface.  It does not assert an attachment to
the punctured $S^6$ family or a compactification of the resulting
threefold over the open parameter disc.
\end{proposition}
\begin{proof}
We first retain the analytic neighbourhood, not merely its normal
line.  At the central parameter the Kato gluing of B.4.10 is
\[
 (x,v)\longmapsto(\kappa a^8xv^8,av).
\]
It is exactly linear in $x$.  Thus a neighbourhood of $E$ is a
neighbourhood of the zero section in the negative line bundle
$N$, with the indicated factor of automorphy.
Its contraction is the cone associated with the section ring
\[
 R(E,N^{-1})=\bigoplus_{m\ge0}H^0(E,N^{-m}).
\]
A line bundle of degree eight on an elliptic curve is very ample
and normally generated.  The cone is normal, its only singularity
is its vertex, and the blow-up of the vertex maximal ideal is
$\operatorname{Tot}(N)$ near its zero section.
One can see the last assertion directly: the degree-one sections
generate the section ring and have no common zero on $E$;
their ratios on the blow-up recover the projective embedding of
$E$, and their common radial coordinate is the coordinate in $N$.
In particular, with $I=\mathcal O_S(-E)$ and $\mathfrak m_p$ the
vertex maximal ideal,
\[
 \mathfrak m_p\mathcal O_S=I.
\]
The local contraction glues to the identity on $S\setminus E$,
giving the stated proper contraction of the compact surface.

Write $C$ for the disjoint rational cycle of the parabolic surface.
The meromorphic two-form
$d\log x\wedge d\log v$ is invariant under the displayed
monomial gluing.  In the toric blow-up charts it has simple poles
along precisely the axis and the rational cycle, and no zeros.
Consequently
\[
 K_S=\mathcal O_S(-E-C).
\]
Also $b_1(S)=1$, by B.4.10.
Put $\mathcal F=\Theta_S(-\log E)$, the locally free sheaf of
vector fields tangent to $E$.
Serre duality and the displayed canonical divisor give
\[
 H^2(S,\mathcal F)^*
 =H^0\bigl(S,\Omega^1_S(\log E)\otimes K_S\bigr)=0.
\]
Indeed there are sheaf injections
\[
 \Omega^1_S(\log E)(-E-C)
       \lhook\joinrel\longrightarrow\Omega^1_S(-C)
       \lhook\joinrel\longrightarrow\Omega^1_S.
\]
Locally the first sends the logarithmic generators
$dx/x,dy$, multiplied by $x$, to the ordinary forms $dx,xdy$.
Finally $H^0(S,\Omega^1_S)=0$ follows already from $b_1(S)=1$.
For completeness, a holomorphic one-form $\alpha$ on a compact
complex surface is closed: Stokes' theorem applied to
$\alpha\wedge\overline{d\alpha}$ makes the integral of the
positive volume form associated with
$d\alpha\wedge\overline{d\alpha}$ vanish.
If $\alpha\ne0$, its real and imaginary parts give independent
real de Rham classes.  An exact real linear combination would
be the differential of a global real pluriharmonic function,
which is constant on the compact surface.  This would give
$b_1(S)\ge2$, a contradiction.

We compute the one possible local higher direct image explicitly.
The restriction of $\mathcal F$ to $E$ is the Atiyah extension
of $N$:
\[
 0\longrightarrow\mathcal O_E
 \longrightarrow\mathcal A_N:=\mathcal F|_E
 \longrightarrow\Theta_E\longrightarrow0.
\]
In normal coordinates related by $x_j=g_{ij}(v)x_i$, the
vertical Euler field is $x_i\partial_{x_i}$, and the changes of
horizontal lifts have cocycle $d\log g_{ij}$.
Thus this extension has class $c_1(N)$ in
$H^1(E,\Omega^1_E)$, which is nonzero because $\deg N=-8$.
Since $\Theta_E\simeq\mathcal O_E$, the connecting map
$H^0(E,\Theta_E)\to H^1(E,\mathcal O_E)$ is an isomorphism.
The same cohomology sequence therefore gives the natural
isomorphism
\[
 H^1(E,\mathcal A_N)\xrightarrow{\ \sim\ }
 H^1(E,\Theta_E)\simeq\C.
\]

The successive quotients for the formal neighbourhood of $E$ are
\[
 I^m\mathcal F/I^{m+1}\mathcal F
     =\mathcal A_N\otimes N^{-m}\qquad(m\ge0).
\]
For $m\ge1$ both line bundles in the induced extension have
degree $8m>0$ on the elliptic curve, so their first cohomology
vanishes.  It follows that all restriction maps
\[
 H^1\bigl(E,\mathcal F/I^{m+1}\mathcal F\bigr)
 \longrightarrow H^1(E,\mathcal F/I^m\mathcal F)
 \quad(m\ge1)
\]
are isomorphisms.  Their common value is
$H^1(E,\mathcal A_N)\simeq\C$.
The analytic theorem on formal functions
\citep[Section~6, Hauptsatz~IIa, p.~59]{Grauert60push}
identifies this inverse limit with the completed stalk of
$R^1\pi_*\mathcal F$ at $p$.
This coherent sheaf is supported at the single point $p$,
so its stalk is a finite-length module and equals its completion.
We have proved
\[
 R^1\pi_*\mathcal F\simeq\C_p,
 \qquad
 H^0(Y,R^1\pi_*\mathcal F)
    \xrightarrow{\ \sim\ }H^1(E,\Theta_E).
\]
The second map is the natural restriction followed by the
Atiyah-extension map just computed, not a dimension count alone.

There is also an equality of sheaves
\[
 \pi_*\mathcal F=\Theta_Y.
\]
To check it at $p$, a holomorphic derivation of the cone
integrates locally to a holomorphic flow: embed the cone in a
polydisc, lift its coefficient functions, and use the ordinary
holomorphic vector-field flow, which preserves the defining ideal.
The flow preserves the unique singular point, hence its maximal
ideal.  Functoriality of the blow-up of that ideal lifts the flow
to the resolution, preserving the exceptional curve.
Its derivative is therefore a section of $\mathcal F$.
Conversely a section of $\mathcal F$ differentiates pullbacks
of holomorphic functions on $Y$ and descends to a derivation,
because $\pi_*\mathcal O_S=\mathcal O_Y$.
Both operations agree off $E$ and are inverse there; uniqueness
extends the equality over $E$.

The global-to-local cohomology map is surjective for a concrete
reason.  Write $a=e^{2\pi i\tau}$ with $\operatorname{Im}\tau>0$,
and vary $\tau$ to $\tau+u$ in the explicit Kato construction,
keeping its translation parameter zero.
The same uniform collar argument as in B.4.10 gives a proper
holomorphic deformation of the pair $(S,E)$.
Its Kodaira--Spencer class belongs to $H^1(S,\mathcal F)$.
Under the natural composite
\[
 H^1(S,\mathcal F)\longrightarrow
 H^0(Y,R^1\pi_*\mathcal F)
 \xrightarrow{\ \sim\ }H^1(E,\Theta_E)
\]
it is the Kodaira--Spencer class of
$\C/(\Z+(\tau+u)\Z)$, which is nonzero.
For example, differentiating its second deck transformation
$z\mapsto z+\tau+u$ gives the usual nonzero period cocycle.
This argument works also at the elliptic curves with extra
automorphisms; it does not differentiate the coarse $j$-invariant.
The target is one dimensional, so the first arrow is onto.

The degree-two part of the Leray sequence is
\[
 H^1(S,\mathcal F)\longrightarrow
 H^0(Y,R^1\pi_*\mathcal F)\longrightarrow
 H^2(Y,\pi_*\mathcal F)\longrightarrow H^2(S,\mathcal F).
\]
The first map is surjective and the last group is zero.
Using $\pi_*\mathcal F=\Theta_Y$ proves
$H^2(Y,\Theta_Y)=0$.

We spell out the local-to-global criterion used in the final
step, since the rational-cycle contraction theorem does not
apply to this elliptic contraction.
For a compact complex space with isolated singularities,
$H^2(\Theta)=0$ makes the restriction from global to local
deformations formally smooth.
Here is the obstruction argument for the present single
singularity.  Take a finite Stein cover with one small
neighbourhood of $p$ and all other members smooth.
For a small extension of Artinian bases, prescribe the lifted
deformation on the singular member and take smooth local
lifts on the other members.  Their pairwise identifications
lift on the smooth Stein intersections.  The failure of the
cocycle identity on triple intersections is a Cech two-cocycle
in $\Theta_Y$ tensored with the square-zero kernel.
Its class vanishes by the proved $H^2$-vanishing; correcting
the identifications by a one-cochain glues the desired lift.
This proves the relative infinitesimal lifting property.

The analytic semiuniversal deformation theory turns this
property into regularity of the induced morphism of analytic
base germs.  The general criterion is
\citet[Lemma~1, p.~93]{Manetti1991D}, stated for reduced compact
complex spaces with isolated singularities and proved by the
local-to-global spectral sequence of the cotangent complex.
It is also used in
\citet[Section~3.2, Proposition~3.3 and the discussion preceding
Proposition~3.4]{DlousskyTeleman2020D}.
There is no hypothesis about the kind of exceptional curve
in this criterion.  A regular morphism of analytic
germs is locally a projection from a product with a smooth
germ.  Consequently a specified analytic disc in the local
deformation base lifts to the global base, without replacing
the prescribed local deformation by a different smoothing
component.
The global family is proper by the compact deformation
construction.  Away from $p$ its central fibre is smooth,
and smoothness is open; compactness permits one common
shrinking of the disc.  A local smoothing at $p$ therefore
gives smooth global noncentral fibres, as asserted.
\end{proof}
